% ============================================================
\chapter{Laplace Equation and Harmonic Functions}
\label{ch:laplace}
% ============================================================

Laplace's equation $\Delta u = 0$ is the queen of elliptic PDEs. Its solutions, \emph{harmonic functions}, model equilibrium states: steady-state temperature, electrostatic potential, irrotational flow. Pierre-Simon Laplace studied it as early as 1782 in the context of Newtonian gravitation, and its properties---maximum principle, infinite regularity, Poisson representation formula---make it one of the best-understood equations in all of mathematical physics.

\section{Introduction and motivation}

Laplace's equation $\Delta u = 0$ models equilibrium states: steady-state
temperature, electrostatic potential in the absence of charges, potential
flow. Its study reveals remarkable regularity, mean value, and maximum
properties.

\begin{definition}[Harmonic function]
A function $u \in C^2(\Omega)$ is \textbf{harmonic} in an open set
$\Omega \subset \R^n$ if:
\begin{equation}\label{eq:laplace}
    \Delta u = \sum_{i=1}^{n} \frac{\partial^2 u}{\partial x_i^2} = 0
    \quad \text{in } \Omega.
\end{equation}
\end{definition}

\begin{definition}[Poisson's equation]
\textbf{Poisson's equation} is the inhomogeneous equation:
\begin{equation}\label{eq:poisson}
    -\Delta u = f \quad \text{in } \Omega.
\end{equation}
\end{definition}

\section{Fundamental solution}

\begin{theorem}[Fundamental solution of the Laplacian]
The \textbf{fundamental solution} of Laplace's equation in dimension $n$ is:
\begin{equation}\label{eq:fundamental_laplace}
    \Gamma(x) = \begin{cases}
        -\frac{1}{2\pi}\ln\abs{x} & \text{if } n = 2,\\[6pt]
        \frac{1}{n(n-2)\omega_n}\frac{1}{\abs{x}^{n-2}} & \text{if } n \geq 3,
    \end{cases}
\end{equation}
where $\omega_n = \frac{\pi^{n/2}}{\Gamma(n/2+1)}$ is the volume of the
unit ball in $\R^n$. It satisfies $-\Delta\Gamma = \delta_0$ in the sense
of distributions.
\end{theorem}

\section{Mean value property}

\begin{theorem}[Mean value property]
If $u$ is harmonic in $\Omega$ and $B(x,r) \subset \Omega$, then:
\begin{equation}\label{eq:mean_value_sphere}
    u(x) = \frac{1}{n\omega_n r^{n-1}}\int_{\partial B(x,r)} u(y)\, dS(y)
    \quad \text{(spherical mean)},
\end{equation}
\begin{equation}\label{eq:mean_value_ball}
    u(x) = \frac{1}{\omega_n r^n}\int_{B(x,r)} u(y)\, dy
    \quad \text{(volume mean)}.
\end{equation}
\end{theorem}

\begin{proof}
Set $\phi(r) = \frac{1}{n\omega_n r^{n-1}}\int_{\partial B(x,r)} u\, dS$.
By the substitution $y = x + r\sigma$, $\sigma \in \partial B(0,1)$:
\[
    \phi(r) = \frac{1}{n\omega_n}\int_{\partial B(0,1)} u(x+r\sigma)\, dS(\sigma).
\]
Differentiating:
\[
    \phi'(r) = \frac{1}{n\omega_n}\int_{\partial B(0,1)}
    \nabla u(x+r\sigma) \cdot \sigma\, dS(\sigma)
    = \frac{1}{n\omega_n r^{n-1}}\int_{\partial B(x,r)}
    \frac{\partial u}{\partial \nu}\, dS.
\]
By the divergence theorem and $\Delta u = 0$:
\[
    \phi'(r) = \frac{1}{n\omega_n r^{n-1}}\int_{B(x,r)} \Delta u\, dy = 0.
\]
So $\phi(r) = \phi(0^+) = u(x)$ by continuity. The volume formula
follows by integrating in $r$.
\end{proof}

\begin{intuition}
The value of a harmonic function at a point equals the average of its
values on any sphere centred at that point. This is the continuous analogue
of the fact that at thermal equilibrium, the temperature at a point is
the average of the surrounding temperatures.
\end{intuition}

\begin{theorem}[Converse of the mean value property]
If $u \in C(\Omega)$ satisfies the mean value property
\eqref{eq:mean_value_sphere} for every ball $B(x,r) \subset \Omega$,
then $u$ is harmonic in $\Omega$ (and in particular $C^\infty$).
\end{theorem}

\section{Maximum principle}

\begin{theorem}[Strong maximum principle for harmonic functions]
\label{thm:strong_max_harmonic}
Let $\Omega \subset \R^n$ be a bounded connected open set and
$u \in C^2(\Omega) \cap C(\overline{\Omega})$ be harmonic in $\Omega$. Then:
\begin{enumerate}
    \item $\max_{\overline{\Omega}} u = \max_{\partial\Omega} u$
          (maximum principle).
    \item If the maximum is attained at an interior point, then $u$ is
          constant in $\Omega$ (strong maximum principle).
\end{enumerate}
The same holds for the minimum.
\end{theorem}

\begin{proof}
Suppose $u$ attains its maximum $M$ at $x_0 \in \Omega$. Let
$r > 0$ be such that $B(x_0, r) \subset \Omega$. By the mean value
property:
\[
    M = u(x_0) = \frac{1}{\omega_n r^n}\int_{B(x_0,r)} u(y)\, dy.
\]
Since $u(y) \leq M$ everywhere and $u(x_0) = M$, the mean can equal $M$
only if $u \equiv M$ in $B(x_0,r)$. The set $\{x : u(x) = M\}$ is
therefore open (and closed in $\Omega$ by continuity). By connectedness,
$u \equiv M$ in $\Omega$.
\end{proof}

\begin{corollary}[Uniqueness of the Dirichlet problem]
The Dirichlet problem $\Delta u = 0$ in $\Omega$, $u = g$ on
$\partial\Omega$, admits at most one solution.
\end{corollary}

\begin{corollary}[Stability]
If $u$ and $v$ are harmonic with $u = g_1$, $v = g_2$ on
$\partial\Omega$, then:
\[
    \max_{\overline{\Omega}} \abs{u - v} \leq \max_{\partial\Omega} \abs{g_1 - g_2}.
\]
\end{corollary}

\section{Regularity of harmonic functions}

\begin{theorem}[$C^\infty$ regularity]
\label{thm:harmonic_regularity}
Every harmonic function is $C^\infty$. More precisely, if $u$ is harmonic
in $\Omega$, then $u \in C^\infty(\Omega)$ and all its partial derivatives
are harmonic.
\end{theorem}

\begin{theorem}[Analyticity]
Every harmonic function is real-analytic: it equals its Taylor series
at every point of its domain.
\end{theorem}

\begin{theorem}[Gradient estimates]
If $u$ is harmonic in $B(x_0, R)$, then:
\begin{equation}\label{eq:gradient_estimate}
    \abs{\nabla u(x_0)} \leq \frac{n}{R}\max_{\partial B(x_0,R)} \abs{u}.
\end{equation}
More generally, for every multi-index $\alpha$:
\[
    \abs{D^\alpha u(x_0)} \leq \frac{C_{n,\abs{\alpha}}}{R^{\abs{\alpha}}}
    \max_{\partial B(x_0,R)} \abs{u}.
\]
\end{theorem}

\section{Harnack's inequality}

\begin{theorem}[Harnack's inequality]
Let $u$ be harmonic and \textbf{nonnegative} in $B(x_0, R) \subset \R^n$.
For any $x \in B(x_0, R)$ with $\abs{x - x_0} = r < R$:
\begin{equation}\label{eq:harnack}
    \frac{R^{n-2}(R-r)}{(R+r)^{n-1}}\, u(x_0) \leq u(x)
    \leq \frac{R^{n-2}(R+r)}{(R-r)^{n-1}}\, u(x_0).
\end{equation}
In particular, on any compact $K \Subset \Omega$, there exists
$C = C(K, \Omega)$ such that:
\[
    \max_K u \leq C \min_K u.
\]
\end{theorem}

\begin{remark}
Harnack's inequality asserts that nonnegative harmonic functions cannot
vary too rapidly. It is fundamental in regularity theory.
\end{remark}

\section{Green's representation formula}

\begin{definition}[Green's function]
The \textbf{Green's function} of the domain $\Omega$ is
$G(x,y) = \Gamma(x-y) - h^y(x)$, where $h^y$ is harmonic in $\Omega$
with $h^y = \Gamma(\cdot - y)$ on $\partial\Omega$.
\end{definition}

\begin{theorem}[Representation formula]
If $u \in C^2(\overline{\Omega})$ is harmonic in $\Omega$, then for
every $x \in \Omega$:
\begin{equation}\label{eq:green_representation}
    u(x) = -\int_{\partial\Omega} u(y)\, \frac{\partial G}{\partial \nu_y}(x,y)\, dS(y).
\end{equation}
For Poisson's equation $-\Delta u = f$:
\begin{equation}\label{eq:poisson_representation}
    u(x) = \int_\Omega G(x,y)\, f(y)\, dy
    - \int_{\partial\Omega} u(y)\, \frac{\partial G}{\partial \nu_y}(x,y)\, dS(y).
\end{equation}
\end{theorem}

\begin{example}[Green's function for the disk]
For the disk $B(0,R) \subset \R^2$, the Green's function is:
\[
    G(x,y) = -\frac{1}{2\pi}\ln\abs{x-y}
    + \frac{1}{2\pi}\ln\!\left(\frac{\abs{x}\abs{y}}{R}
    \abs{\frac{x}{\abs{x}^2} \cdot R^2 - y}\right).
\]
The representation formula recovers the \textbf{Poisson kernel}.
\end{example}

\section{Liouville's theorem}

\begin{theorem}[Liouville's theorem]
Every bounded harmonic function on $\R^n$ is constant.
\end{theorem}

\begin{proof}
Let $u$ be bounded harmonic on $\R^n$, $\abs{u} \leq M$. For any
$x_0 \in \R^n$ and $R > 0$, the gradient estimate gives:
\[
    \abs{\nabla u(x_0)} \leq \frac{n}{R}\, M.
\]
Letting $R \to \infty$, we obtain $\nabla u(x_0) = 0$ for all $x_0$,
so $u$ is constant.
\end{proof}

\section{Exercises}

\begin{exercise}
Show that if $u$ is harmonic in $B(0,R) \setminus \{0\} \subset \R^2$
and bounded, then the singularity at $0$ is removable (removable
singularity theorem).
\end{exercise}

\begin{exercise}
Use the Poisson formula to solve the Dirichlet problem in the disk
$B(0,1)$ with $u(1,\theta) = \cos(2\theta) + 3\sin(\theta)$.
\end{exercise}

\begin{exercise}
Show that if $u$ is harmonic and nonnegative in $\R^n$, then $u$ is
constant (consequence of Harnack).
\end{exercise}

\begin{exercise}
Compute the Green's function of the half-space $\{x_n > 0\}$ by the
method of images.
\end{exercise}

\begin{exercise}
Show that the mean value of a harmonic function on a sphere equals its
value at the centre. What happens for subharmonic functions
($\Delta u \geq 0$)?
\end{exercise}

\begin{exercise}
Let $u$ be harmonic in $\Omega = B(0,1) \setminus \overline{B(0,1/2)}
\subset \R^3$, with $u = 0$ on $\abs{x} = 1/2$ and $u = 1$ on $\abs{x} = 1$.
Find $u$ in the form $u(r) = a + b/r$.
\end{exercise}

\begin{exercise}[Schwarz reflection principle]
Let $u$ be harmonic in the half-disk $\{r < 1, 0 < \theta < \pi\}$
with $u = 0$ on the diameter $[-1,1]$. Show that $u$ extends to a
harmonic function on the full disk by odd reflection.
\end{exercise}
